\section{Conclusion et Perspectives}

En conclusion, ce stage aura permis de mettre en place une méthodologie de la 
collecte des données jusqu'à la prédiction et l'évaluation des modèles. 

Les résultats obtenus indiquent clairement que les modèles de Machine Learning 
surpassent l'approche statistique classique, tant en termes de capacité prédictive 
que de rentabilité. Ces performances mettent toutefois en évidence un élément crucial : 
la rentabilité à long terme dépend directement de la sélectivité du modèle. 
Ce dernier doit restreindre son volume de paris via un seuil optimal pour ne pas 
prendre de risque inutile et perdre de l'argent.

Plusieurs pistes restent néanmoins ouvertes pour approfondir ce 
travail. Tout d'abord, l'extension de la base de données constituerait une amélioration 
naturelle : intégrer davantage de championnats (compétitions européennes secondaires, 
voire d'autres confédérations) et un historique plus long permettrait d'augmenter la 
taille de l'échantillon d'entraînement et d'évaluer la généralisation des modèles à des 
contextes footballistiques différents. De même, l'intégration de données en direct, telles 
que les statistiques de possession, les changements tactiques, les blessures des joueurs ou 
la composition des équipes, pourrait significativement enrichir l'information disponible 
et ouvrir la voie à une prédiction dynamique, actualisée tout au long de la rencontre.

Toutefois, il reste relativement compliqué au vu de la visualisation section \ref{sec:visu} 
de prédire avec certitude les matchs de foot. En effet, grace aux méthodes de visualisation 
nous voyons clairement que les données sont relativement peu séparables, ce qui explique 
nos résultats plutôt faible en terme de performances prédictive et financière. 
La méthodologie de sélection de variables pourrait elle aussi être affinée, par 
exemple en explorant des approches d'embedding ou des méthodes de réduction de dimension 
non linéaires en amont de la sélection, afin de capturer des interactions plus complexes 
entre les descripteurs. 